%! Setting Up a Test for a Population Mean — Unit 7, Lesson 7.4 — Group Activity (Answer Key)
\documentclass[10pt]{article}
\usepackage{apstats-article}
\usepackage{apstats-key}   % pulls in -boxes; do NOT also load -boxes
\newcommand{\TallMath}[1]{$\displaystyle #1\rule[-1.4em]{0pt}{3.2em}$}

\begin{document}

\pageheader{Unit 7, Lesson 7.4}{Group Activity --- Answer Key}
\namepartnerperiod

\vspace{0.15in}

\textit{For each scenario, (a) name the appropriate inference method, (b) state $H_0$ and
$H_a$ (including the direction of $H_a$), and (c) verify both conditions. Do \textbf{not}
compute a test statistic yet.}

\begin{scenariobox}[Scenario A --- Sleep Duration]{sky}
\textbf{Tier A \& R.} \quad
The American Academy of Sleep Medicine recommends that teenagers get at least 8 hours of sleep
per night. A school counselor believes students at her school sleep \emph{less} than 8 hours
on average. She randomly selects 20 students and records their nightly sleep duration (hours).
The dotplot shows the distribution is roughly symmetric with no outliers.

\begin{enumerate}[label=\alph*.]
  \item Inference method: \ans{one-sample $t$-test for $\mu$} \quad $df = $ \ans{19}
  \item $H_0$: \ans{$\mu = 8$} \quad $H_a$: \ans{$\mu < 8$}

        Direction of $H_a$: \ans{left (one-sided)} \quad Justify: \ansline{The counselor believes students sleep \emph{less than} 8 hours --- not just different.}

  \item Conditions:
        \begin{itemize}
          \item Random: \ansline{Random sample of 20 students --- met.}
          \item Normal ($n = 20 < 30$): \ansline{Dotplot shows roughly symmetric, no outliers --- Normal condition met.}
        \end{itemize}
\end{enumerate}
\end{scenariobox}

\begin{scenariobox}[Scenario B --- Package Weights]{sky}
\textbf{Tier A \& R.} \quad
A snack company labels its chip bags ``Net weight: 1.5 oz.'' A quality-control inspector
suspects the \emph{true mean} weight differs from 1.5 oz. She randomly selects 40 bags and
weighs each one. The histogram of weights appears approximately bell-shaped.

\begin{enumerate}[label=\alph*.]
  \item Inference method: \ans{one-sample $t$-test for $\mu$} \quad $df = $ \ans{39}
  \item $H_0$: \ans{$\mu = 1.5$} \quad $H_a$: \ans{$\mu \ne 1.5$}

        Direction of $H_a$: \ans{two-sided} \quad Justify: \ansline{The question is whether the weight \emph{differs} --- no specific direction stated.}

  \item Conditions:
        \begin{itemize}
          \item Random: \ansline{Random sample of 40 bags --- met.}
          \item Normal ($n = 40 \ge 30$): \ansline{$n = 40 \ge 30$; CLT applies --- met.}
        \end{itemize}
\end{enumerate}
\end{scenariobox}

\begin{scenariobox}[Scenario C --- Matched Pairs: Reading Fluency]{sky}
\textbf{Tier A \& E.} \quad
A reading specialist administers a fluency test to 12 students at the beginning and end of a
semester-long intervention. She wants to determine whether the intervention \emph{improved}
reading speed (words per minute). The differences (end $-$ begin) are available for analysis.
A boxplot of the differences shows no strong skewness or outliers.

\begin{enumerate}[label=\alph*.]
  \item What type of data is this? Circle: \quad one sample \quad \ans{matched pairs}

  \item Define the parameter $\mu_d$ in context: \ansline{The true mean difference in reading speed (end $-$ begin) for all students in this intervention program.}

  \item Define the order of subtraction and state hypotheses:

        $d_i = $ \ans{end} $-$ \ans{begin} \qquad $H_0$: \ans{$\mu_d = 0$} \quad $H_a$: \ans{$\mu_d > 0$}

  \item Verify conditions:
        \begin{itemize}
          \item Random/Independence: \ansline{Students were treated as a convenience sample from the intervention group --- discuss whether this meets the random condition; note any limitation.}
          \item Normal ($n = 12 < 30$): \ansline{Boxplot of differences shows no strong skewness or outliers --- Normal condition met.}
        \end{itemize}
\end{enumerate}
\end{scenariobox}

\begin{reflectionbox}
Scenarios A and B both involve unknown $\sigma$. What is the key difference in the Normal
condition check between them? Why does sample size matter here?
\ansline{Scenario A: $n = 20 < 30$ --- must inspect the sample distribution (dotplot/boxplot) for strong skewness or outliers.}
\ansline{Scenario B: $n = 40 \ge 30$ --- CLT guarantees approximate normality regardless of population shape; no graph check needed.}
\end{reflectionbox}

\begin{teachernote}
Scenario C is intentionally ambiguous about randomness. Use it to discuss the difference between
a convenience sample and a random sample, and how to note limitations in the conclusion.
\end{teachernote}

\end{document}
